% Section 1.3: Searching for Dark Matter
% Structure: Detection triangle overview → Direct detection → Collider searches → Why indirect detection

\section{Searching for Dark Matter}
\label{sec:1.3}

With the theoretical target established --- a particle in the GeV--TeV mass range whose annihilation cross-section is anchored to the thermal relic value --- the question becomes how to detect it experimentally.
The answer is not a single experiment but a coordinated campaign that tackles the problem from three distinct directions.
If dark matter interacts with ordinary matter through any non-gravitational force, then the same coupling that governs annihilation in the early universe also mediates scattering in underground laboratories and production at particle colliders.
These three approaches --- direct detection, collider production, and indirect detection --- form what is commonly known as the ``detection triangle''.

Each vertex of this triangle probes a different kinematic regime of the same underlying interaction.
Direct and indirect detection both observe dark matter in its natural habitat: the galactic halo, where particles drift at non-relativistic velocities of $v \sim 10^{-3}\,c$.
Colliders, by contrast, collide particles at energies far exceeding the dark matter mass, accessing a fully relativistic regime.
Because the same dark matter--Standard Model coupling can yield very different cross-sections in the non-relativistic regime probed by direct and indirect detection and the relativistic regime probed at colliders --- sometimes suppressed by powers of the dark matter velocity or the momentum transfer --- an interaction that produces a strong signal in one channel can be entirely invisible in another.
A signal in one corner combined with a non-detection in another would therefore narrow the type of interaction at play, providing complementary information that no single channel could yield.
A signal in all three would constitute an unambiguous identification.

\begin{figure}[t]
	\centering
	\includegraphics[width=0.95\textwidth]{figures/DMdetections.pdf}
	\caption{The detection triangle: three complementary strategies for detecting dark matter, each corresponding to a different reading of the same underlying interaction vertex.
		\textit{Left}: direct detection --- dark matter scatters off a Standard Model target in a laboratory.
		\textit{Center}: indirect detection --- dark matter annihilates into Standard Model particles in astrophysical environments.
		\textit{Right}: collider production --- Standard Model particles produce dark matter at accelerators.
		Figure from Cirelli et al.~\cite{Cirelli:2024ssz}.
	}
	\label{fig:detection_triangle}
\end{figure}

This section reviews each strategy in turn, with emphasis on the physical principles rather than experimental details.
We close with the argument for why indirect detection via gamma-rays --- the approach pursued throughout this thesis --- provides the most direct test of the WIMP hypothesis.

\subsection{Direct Detection}
\label{sec:1.3.1}

Direct detection experiments search for the faint recoil produced when a dark matter particle from the galactic halo scatters off an atomic nucleus inside a carefully shielded underground detector.
The kinematics are straightforward: a WIMP with mass $M_\chi$ and velocity $v \sim 10^{-3}\,c$ strikes a target nucleus of mass $m_A$ essentially at rest, transferring a recoil energy $E_R \sim \mu_A^2 v^2 / m_A$, where $\mu_A = M_\chi m_A / (M_\chi + m_A)$ is the reduced mass of the system~\cite{Dodelson:2020bqr}.
For a WIMP and a xenon nucleus of comparable mass ($M_\chi \sim m_A \sim 100$~GeV), this yields $E_R \sim 20$~keV --- small enough that detecting it requires extraordinarily sensitive instruments, but large enough to be within reach of modern technology.
In addition to the recoil energy spectrum, the Earth's orbital motion around the Sun modulates the WIMP flux by ${\sim}\,5\%$ annually (peaking around June~2), providing a distinctive time-dependent signature~\cite{Drukier:1986tm}.
The DAMA/LIBRA experiment has observed a persistent annual modulation signal in NaI(Tl) crystals over more than two decades~\cite{Bernabei:2018yyw}, but its interpretation as WIMPs is in strong tension with null results from other experiments using different target materials, and the claimed signal remains controversial~\cite{ANAIS-112:2025fne}.

The physics of the scattering divides into two categories determined by the non-relativistic limit of the interaction~\cite{Cirelli:2024ssz,Hooper:2024}.
In \textbf{spin-independent (SI)} scattering, the dark matter couples to the nucleon number density without reference to nuclear spin.
In the limit of small momentum transfer --- where the de~Broglie wavelength of the dark matter particle exceeds the nuclear radius --- the amplitudes from individual nucleons add coherently, producing a total cross-section that scales as $\sigma_\text{SI} \propto A^2$, where $A$ is the atomic mass number~\cite{Cirelli:2024ssz,Hooper:2024}.
This quadratic enhancement is the primary reason why the most sensitive experiments use heavy target elements, particularly xenon ($A \approx 131$).
\textbf{Spin-dependent (SD)} scattering, by contrast, couples to the nuclear spin $J$ and receives no coherent enhancement, leading to bounds that are typically three to five orders of magnitude weaker than the SI case~\cite{Cirelli:2024ssz}.

The current generation of direct detection experiments is dominated by dual-phase liquid xenon time projection chambers (TPCs), which detect both the scintillation light (S1~signal) and the ionization charge (S2~signal) from a nuclear recoil event~\cite{Cirelli:2024ssz}.
The ratio S2/S1 provides powerful discrimination between nuclear recoils (the expected dark matter signal) and electron recoils (the dominant background).
Three experiments lead the field: \textbf{LZ} (LUX-ZEPLIN)~\cite{LZ:2022lsv} in the United States, \textbf{XENONnT}~\cite{XENON:2023cxc} at the Gran Sasso laboratory in Italy, and \textbf{PandaX-4T}~\cite{PandaX-4T:2021bab} in China.
Together, they have pushed the upper limit on the SI WIMP-nucleon cross-section down to $\sigma_\text{SI} \lesssim 10^{-47}$~cm$^2$ near the minimum of the exclusion curve at $M_\chi \sim 30$~GeV~\cite{LZ:2022lsv,XENON:2023cxc,PandaX-4T:2021bab}.
At higher masses ($M_\chi \sim 100$~GeV), the limits relax to ${\sim}\,10^{-46}$~cm$^2$ as the local number density of dark matter particles decreases~\cite{LZ:2024zvo}.
Beyond these noble liquid TPCs, a complementary program of cryogenic bolometers (SuperCDMS, CRESST-III) and semiconductor-based detectors (SENSEI, DAMIC) extends sensitivity to sub-GeV dark matter by searching for electron recoils rather than nuclear recoils~\cite{Cirelli:2024ssz}.

Standard Model neutrinos --- solar, atmospheric, and from the diffuse supernova background --- scatter coherently off nuclei and produce recoils kinematically indistinguishable from a WIMP signal; this irreducible background is the \textbf{neutrino fog} (or neutrino floor).
The question of whether it would ever actually show up in data was settled in 2021, when XENON1T recorded the first observation of solar $^8$B coherent elastic neutrino-nucleus scattering (CE$\nu$NS) in a dark matter detector~\cite{XENON:2020gfr}.
PandaX-4T and XENONnT have since reported independent indications at roughly $2.6\sigma$ and $2.7\sigma$, consistent with the expected solar rate~\cite{PandaX:2024muv,XENON:2024hup}; XENONnT also ran the first dedicated dark matter search inside the fog~\cite{XENON:2024hup}, and LZ set the world-best spin-independent WIMP bound while reporting its own $^8$B indications~\cite{LZ:2024zvo}.
What this means practically is that the fog is already a feature of current data, not a future cliff: sensitivity gains slow from here, scaling as the inverse square root of accumulated exposure rather than inversely~\cite{OHare:2021utq}, and the next generation of detectors --- notably the DARWIN/XLZD program~\cite{DARWIN:2016hyl} --- will reach multi-ton xenon exposures designed to operate within the fog rather than to avoid it.
Pushing substantially beyond it will likely require directional detectors capable of resolving the apparent ``dark matter wind'' from the Cygnus constellation against the isotropic neutrino background.

Direct detection excels in the mass range where the WIMP mass is comparable to the target nucleus mass --- roughly tens of GeV for xenon targets~\cite{Cirelli:2024ssz}.
At lower masses, the recoil energy drops below experimental thresholds; at higher masses, the local number density decreases as $n \propto 1/M_\chi$, reducing the event rate.
These mass-dependent limitations, combined with our entry into the neutrino fog, motivate a complementary search strategy.
Direct detection constrains the WIMP--nucleon scattering cross-section but is insensitive to the annihilation cross-section $\langle \sigma v_\mathrm{rel} \rangle$ that determines the thermal relic abundance --- the observable probed by indirect detection.

We focus on WIMP--nucleus scattering as a methodological choice rather than as the full scope of direct detection, since other well-motivated candidates such as the axions introduced in Section~\ref{sec:1.2.1} require entirely different experimental strategies.
The standard probe is a resonant-cavity \emph{haloscope}: in a strong static magnetic field, an axion from the local galactic halo can convert into a microwave photon inside a high-quality-factor cavity tuned to the axion mass, producing a narrow excess above the thermal noise floor.
Searches of this type include ADMX~\cite{ADMX:2019uok} and HAYSTAC~\cite{HAYSTAC:2024jch}, both of which probe the QCD axion parameter space.
Since this thesis concerns the WIMP scenario and its gamma-ray indirect signatures, we do not pursue axion direct detection further, and refer the reader to the Snowmass community report~\cite{Adams:2022pbo} for a survey of the axion experimental program.

For comprehensive recent treatments of direct detection covering experimental techniques, current bounds, and the path beyond the neutrino fog, we refer the reader to the reviews by Schumann~\cite{Schumann:2019eaa} and Cirelli et al.~\cite{Cirelli:2024ssz}, and to Hooper's book~\cite{Hooper:2024}.

\subsection{Collider Searches}
\label{sec:1.3.2}

If dark matter couples to quarks or gluons, it can be pair-produced in the high-energy collisions at the Large Hadron Collider.
The challenge is detection: dark matter particles carry no electric or color charge and escape the detector without leaving a trace, just as neutrinos do~\cite{Cirelli:2024ssz,Hooper:2024}.
The only unavoidable signature is \textbf{missing transverse energy} ($E_T^\text{miss}$), inferred from the momentum imbalance of the visible particles in the event~\cite{Cirelli:2024ssz}.

Pure pair-production of dark matter ($pp \to \chi\chi$) produces no visible particles and goes entirely unnoticed.
A detectable signal requires that an additional Standard Model particle --- a jet, a photon, or a $W/Z$ boson --- is radiated from the initial state, recoiling against the invisible dark matter pair.
These are the \textbf{mono-X searches}: mono-jet, mono-photon, mono-$Z$, and mono-Higgs, each paying a cross-section penalty of roughly $g^2/(4\pi)^2 \sim 10^{-3}$ for the additional emission~\cite{Cirelli:2024ssz}.
A complementary channel probes \textbf{invisible Higgs decays}: if the dark matter is lighter than half the Higgs boson mass ($M_\chi < M_h/2$), the Higgs can decay directly into a dark matter pair.
Current LHC measurements constrain the invisible branching ratio to $\text{BR}_\text{inv} \lesssim 0.11$ at 95\%~C.L., essentially excluding Higgs-portal thermal relic dark matter below $M_h/2 \approx 63$~GeV~\cite{ATLAS:2023tkt,Cirelli:2024ssz}.
The earliest collider analyses treated dark matter production as a contact interaction between dark matter and quarks, a model-independent framework that breaks down at the LHC, where the partonic collision energy routinely approaches or exceeds the mediator mass and the effective-operator expansion is no longer self-consistent.
We therefore work with \textbf{simplified models} that keep an explicit dark matter particle, an explicit mediator (most commonly an $s$-channel vector $Z'$ or a pseudoscalar), and explicit couplings to quarks and to dark matter, which lets us confront the same parameter space with direct and indirect detection bounds~\cite{Boveia:2018yeb}.
One consequence is worth spelling out: any mediator coupled to quarks strongly enough to be produced at the LHC must also decay back into quark-antiquark pairs, and the absence of a di-jet resonance is often a stronger constraint than the corresponding missing-energy search~\cite{Boveia:2018yeb}.
Combined with mono-jet limits, current LHC data exclude $s$-channel mediators across roughly the $0.1$--$4$~TeV range unless the quark coupling is much smaller than the dark matter coupling, $g_q \ll g_\text{DM}$.
For colored mediators such as squarks and gluinos in supersymmetric realizations, the larger QCD production rates push the lower bound on their masses to roughly $1$--$2$~TeV~\cite{Cirelli:2024ssz,Hooper:2024}.
These collider bounds probe interaction strengths that direct detection cannot reach, because the production cross-section at LHC energies does not feel the velocity and momentum suppressions that control non-relativistic dark matter-nucleus scattering.

Despite their power, collider experiments face a fundamental limitation: even if a new invisible particle is produced and characterized, the LHC cannot establish that it is cosmologically stable or that it constitutes the dark matter observed in gravitational systems.
Confirming the cosmological connection --- that the particle froze out in the early universe with the correct relic abundance --- requires measuring the annihilation cross-section directly.
This is the domain of indirect detection.

For a comprehensive recent treatment of dark matter searches at colliders, including the simplified-model framework and the LHC mono-X programme, we refer the reader to the review by Boveia and Doglioni~\cite{Boveia:2018yeb}, and to the broader reviews by Cirelli et al.~\cite{Cirelli:2024ssz} and Hooper~\cite{Hooper:2024}.

\subsection{Why Indirect Detection?}
\label{sec:1.3.3}

Among the three strategies, indirect detection occupies a unique position: it is the only approach that directly probes the annihilation cross-section $\langle \sigma v_\text{rel} \rangle$ --- the same quantity that determines the thermal relic abundance~\cite{Hooper:2024}.
For a standard $s$-wave WIMP, the cross-section at freeze-out is essentially the same as the one operating in today's galactic halos (see Section~\ref{sec:1.2.4}), providing a firm, parameter-free experimental target~\cite{Steigman:2012nb,Hooper:2024}.
Neither direct detection nor collider production offers this direct connection to cosmology.

Dark matter annihilation (or decay) in astrophysical environments produces all stable Standard Model particles --- gamma-rays, neutrinos, electrons, positrons, antiprotons, and heavier anti-nuclei --- each carrying information about the mass, coupling, and spatial distribution of the dark sector~\cite{Cirelli:2010xx,Arina:2023eic,Cirelli:2024ssz}.
Among these messengers, \textbf{gamma-rays} occupy a privileged position.
Photons are electrically neutral and travel in straight lines, preserving both the angular distribution and the spectral signature of their source~\cite{Cirelli:2024ssz,Pinetti:2021jjs}, whereas charged cosmic rays --- positrons, antiprotons --- are deflected by galactic magnetic fields, erasing all directional information and distorting the energy spectrum through diffusion and energy losses~\cite{Cirelli:2024ssz}.
The universe is also largely transparent to gamma-rays up to TeV-scale energies over galactic distances, so they reach the detector without significant absorption~\cite{Hooper:2024}.
Neutrinos share the advantage of being neutral and undeflected, but their interaction cross-sections are so small that detecting them requires cubic-kilometer-scale instruments and still yields far fewer events than a gamma-ray telescope of comparable collecting area~\cite{Cirelli:2024ssz}.
Gamma-rays therefore provide one of the cleanest windows on dark matter annihilation: they carry uncontaminated morphological information (where the dark matter is concentrated) and spectral information that informs the dark matter mass and annihilation channel, modulo the degeneracies arising from spectral similarity between final states and from astrophysical uncertainties.

The potential of gamma-ray indirect detection is underscored by the most debated anomaly in the field: the \textbf{Galactic Center excess} (GCE), a statistically significant excess of GeV-scale gamma-rays from the inner Galaxy whose spectrum and morphology are consistent with dark matter annihilation but also with an unresolved population of millisecond pulsars~\cite{Goodenough:2009gk,Hooper:2010mq,Daylan:2014rsa,Macias:2019omb}.
We discuss the GCE in detail in Section~\ref{sec:targets}, after developing the density profile and $J$-factor formalism needed to interpret it quantitatively; the dedicated analysis is presented in Chapter~\ref{ch:4}.

To circumvent the significant backgrounds of the Galactic Center, astronomers have turned to \textbf{dwarf spheroidal galaxies} (dSphs) --- satellite galaxies of the Milky Way that are among the most dark-matter-dominated objects known, with mass-to-light ratios reaching into the thousands~\cite{Fermi-LAT:2015att}.
These systems are essentially devoid of gas, active star formation, and pulsars, producing negligible astrophysical gamma-ray emission.
While the expected dark matter flux from any individual dwarf is orders of magnitude smaller than from the Galactic Center, the absence of backgrounds makes them extraordinarily clean targets~\cite{Hooper:2024}.
By stacking observations of dozens of these systems, the Fermi-LAT collaboration has derived some of the tightest constraints on the annihilation cross-section to date, excluding the thermal relic value for masses below ${\sim}\,100$~GeV in the $b\bar{b}$ channel~\cite{Fermi-LAT:2015att,Hooper:2024}.
These constraints carry systematic uncertainties related to the dark matter density profiles inferred from stellar kinematics~\cite{Bonnivard:2015xpq} --- an issue we revisit in Chapter~\ref{ch:5}.

For comprehensive recent treatments of indirect detection covering target classes, propagation effects, and the multi-messenger landscape, we refer the reader to the reviews by Cirelli et al.~\cite{Cirelli:2024ssz} and Hooper~\cite{Hooper:2024}.

With the case for gamma-ray indirect detection established, the next section develops the underlying formalism.
We separate the predicted flux into a particle-physics factor and an astrophysical $J$-factor, and survey the catalogue of observational targets that together turn this strategy into a quantitative search.
